\section{Lista de implementación: empresas, equipos de investigación y quienes formulan políticas}

\subsection{Para quien aplica esto en la empresa: reforzar primero el «proceso verificable»}

Para la mayoría de las empresas, el orden correcto no es «perseguir primero el modelo más reciente», sino «definir primero una tarea verificable». En la entrevista se menciona varias veces que el código y las matemáticas avanzan rápido porque cuentan con un verificador de aceptación claro (tests / checker).

\begin{importantbox}{Ciclo mínimo de 6 pasos para la implementación empresarial}
\begin{enumerate}
    \item Elegir 1 o 2 procesos de alto valor con aceptación automatizable (por ejemplo, corrección de código, extracción de documentos, enrutamiento de tickets)
    \item Establecer un conjunto de evaluación unificado, con un modelo base y un Prompt base fijos
    \item Adoptar Pass@k más el costo por éxito como métricas centrales
    \item Registrar de manera obligatoria las trayectorias de fallo (entrada, llamadas a herramientas, tipo de error, ruta de recuperación)
    \item Diseñar puntos de intervención humana y un mecanismo de reversión
    \item Revisar semanalmente el Top 10 de fallos antes de decidir si se cambia de modelo
\end{enumerate}

Si desde el inicio se persigue un «Agente totalmente autónomo», por lo general se cae en una doble trampa de confiabilidad y gobernanza.
\end{importantbox}

\subsection{Para los equipos de modelos fundacionales: convertir la «ventaja de entrenamiento» en «ventaja de servicio»}

La entrevista señala una fractura frecuente: el modelo va a la delantera en artículos y benchmarks, pero la concurrencia, la latencia y la disponibilidad del lado del servicio no la siguen, y al final el valor lo capta una plataforma de alojamiento externa. Para evitar esta fractura, hay que optimizar al menos tres cosas de manera simultánea:
\begin{itemize}
    \item \textbf{Arquitectura de inferencia}: estrategia de KV Cache, estrategia de procesamiento por lotes, programación de contexto largo
    \item \textbf{Contrato de producto}: qué capacidades se pueden prometer y cuáles solo se abren de forma experimental
    \item \textbf{Experiencia del desarrollador}: interpretabilidad de errores, transparencia de cuotas, ritmo de las actualizaciones de versión
\end{itemize}

\begin{table}[H]
\centering
\begin{tabular}{lp{8.5cm}}
\toprule
\textbf{Objetivo} & \textbf{Acción recomendada} \\
\midrule
Reducir el costo de las alucinaciones & Exigir el uso de herramientas más un checker en las rutas verificables, y reducir los ciclos cerrados que dependen solo del lenguaje natural \\
Reducir el costo de los reintentos & Emitir códigos de error estructurados y las causas del fallo, para reducir los reintentos a ciegas \\
Elevar la eficiencia de la migración & Al lanzar un nuevo modelo, ofrecer una capa de compatibilidad y un interruptor de reversión, para que el desarrollador no tenga que reescribir todo de una vez \\
\bottomrule
\end{tabular}
\caption{Acciones clave para pasar de «liderazgo del modelo» a «liderazgo de la plataforma»}
\end{table}

\subsection{Para quienes formulan políticas: convertir los objetivos regulatorios en indicadores medibles}

La lección más importante del programa sobre los controles de exportación es esta: si el objetivo se limita a «restringir los flops de entrenamiento», en el nivel de ejecución surgirán numerosas rutas alternativas (memoria, interconexión, arrendamiento, entidades desdobladas). La política necesita pasar de la «restricción nominal» al «resultado medible».

Un marco de política más ejecutable podría incluir:
\begin{itemize}
    \item Indicadores de capacidad de entrenamiento: umbral de acceso a cómputo, disponibilidad de componentes clave
    \item Indicadores de capacidad de despliegue: capacidad de operar de manera sostenida clústeres de inferencia a gran escala
    \item Indicadores de capacidad de difusión: transparencia de las cadenas de arrendamiento y reventa transfronterizas
    \item Indicadores de capacidad de riesgo: resiliencia de la infraestructura crítica (red eléctrica, telecomunicaciones) frente a ataques con IA
\end{itemize}

\begin{warningbox}{No basta con vigilar «si entraron o no los chips»; también hay que vigilar «si se formó o no la capacidad»}
Si la evaluación de la política solo observa el flujo del hardware, puede resultar estadísticamente «en cumplimiento» y, en términos de capacidad, «desprotegida». Los casos reales que aporta la entrevista (arrendamiento distribuido, transbordo en zonas grises) muestran que los indicadores de la política deben vincularse con la capacidad finalmente formada; de lo contrario, es fácil caer en una falsa tranquilidad.
\end{warningbox}

\subsection{Resumen de la sección}

Ya se trate de empresas, equipos de modelos o quienes formulan políticas, lo importante es convertir el objetivo en un ciclo cerrado verificable, iterable y auditable. La competencia en IA no es un enfrentamiento de una sola vez, sino un proceso continuo de convertir la ventaja técnica en una ventaja de ejecución estable.
